\documentclass[aspectratio=169, 11pt, xcolor={table}]{beamer}
% Shared style for the course slides, based on presentation_text/main.tex.
% Dark blue removed: every slide (cover included) uses the light background
% of the template's content slides, with teal + orange as accents.

\usepackage{amsmath}
\usepackage{amssymb}
\usepackage{booktabs}
\usepackage{colortbl}
\usepackage{tabularx}
\usepackage{tikz}
\usepackage[skins]{tcolorbox}
\usepackage{fontawesome5}
\usepackage{polyglossia}

\usetikzlibrary{positioning, calc, arrows.meta, shapes.geometric}

%------------------------------------------------
% Languages and fonts (compile from slides/ with XeLaTeX or tectonic)
%------------------------------------------------
\setdefaultlanguage[numerals=maghrib]{arabic}
\setotherlanguage{english}
\usefonttheme{serif}

\newfontfamily\arabicfont[
    Script=Arabic, Scale=0.98, Path=../fonts/, Extension=.ttf,
    UprightFont=*-Regular, BoldFont=*-Bold,
    ItalicFont=*-Italic, SlantedFont=*-Italic,
    BoldItalicFont=*-BoldItalic, BoldSlantedFont=*-BoldItalic]{Amiri}
\newfontfamily\arabicfontsf[
    Script=Arabic, Scale=0.98, Path=../fonts/, Extension=.ttf,
    UprightFont=*-Regular, BoldFont=*-Bold,
    ItalicFont=*-Italic, SlantedFont=*-Italic,
    BoldItalicFont=*-BoldItalic, BoldSlantedFont=*-BoldItalic]{Amiri}
% Heading font
\newfontfamily\khatt[Script=Arabic, Scale=1.05, Path=../fonts/,
    BoldFont=ae_AlBattar.ttf]{ae_AlBattar.ttf}

%------------------------------------------------
% Colors (template palette without the navy tones)
%------------------------------------------------
\definecolor{brandTeal}{HTML}{00B4D8}
\definecolor{brandOrange}{HTML}{FF7043}
\definecolor{brandLight}{HTML}{F8F9FA}
\definecolor{brandDark}{HTML}{263238}
\definecolor{brandMuted}{HTML}{6C757D}
\colorlet{tealDeep}{brandTeal!70!black}

\setbeamercolor{background canvas}{bg=brandLight}
\setbeamercolor{normal text}{fg=brandDark}
\setbeamercolor{structure}{fg=brandTeal}
\setbeamertemplate{navigation symbols}{}

% beamer.cls sets \raggedright globally; under RTL that centers short lines
% or pushes them left, so go back to justified text (flush right).
\rightskip=0pt\relax
% Lists call \raggedright too; make it a no-op for RTL slides.
\AtBeginDocument{\renewcommand{\raggedright}{\rightskip=0pt\relax}}

%------------------------------------------------
% Frame title, footer, lists
%------------------------------------------------
\setbeamertemplate{frametitle}{%
    \vspace{0.35cm}
    \noindent{\khatt\bfseries\Large\color{brandDark}\insertframetitle}
    \par
    \vskip-2pt
    \noindent{\color{brandTeal}\rule{\textwidth}{1.5pt}}
    \vspace{0.05cm}
}

% Footer: page number only
\setbeamertemplate{footline}{%
    \hbox to \paperwidth{\hfill
        {\color{brandMuted}\scriptsize\LR{\insertframenumber\,/\,\inserttotalframenumber}}%
        \hspace{0.6cm}}%
    \vspace{0.25cm}}

\setbeamertemplate{itemize item}{\color{brandTeal}\small\faChevronLeft}
\setbeamertemplate{itemize subitem}{\color{brandMuted}\small\faAngleLeft}
\setbeamertemplate{enumerate item}{\color{tealDeep}\bfseries\insertenumlabel)}

% Tables: every cell framed by a visible border
\newcolumntype{R}{>{\raggedleft\arraybackslash}X}
\newcolumntype{C}{>{\centering\arraybackslash}X}
\newcommand{\ras}[1]{\cellcolor{brandTeal!20}\textbf{#1}}
% Colored text inside a table cell (\leavevmode keeps it on the row baseline)
\newcommand{\cell}[2]{\leavevmode\textcolor{#1}{#2}}
\renewcommand{\arraystretch}{1.35}
\setlength{\arrayrulewidth}{0.6pt}
\arrayrulecolor{brandMuted!70}

%------------------------------------------------
% Boxes (the two boxes of the template's 2nd slide)
%------------------------------------------------
\tcbset{
    brandbox/.style={
        enhanced, arc=4pt, outer arc=4pt, boxrule=1.5pt,
        coltitle=white, coltext=brandDark, fonttitle=\khatt\bfseries,
        toptitle=0.8mm, bottomtitle=0.8mm,
        top=1.2mm, bottom=1.2mm, left=3mm, right=3mm,
        before skip=4pt, after skip=4pt}
}
% Teal box: definitions, objectives, summaries
\newtcolorbox{tbox}[2][]{brandbox, colback=brandTeal!10,
    colframe=brandTeal, title={#2}, #1}
% Orange box: warnings
\newtcolorbox{obox}[2][]{brandbox, colback=brandOrange!10,
    colframe=brandOrange, title={#2}, #1}
% Card (template's 3rd slide, in teal)
\newtcolorbox{card}[2][]{enhanced, colback=brandTeal!6, colframe=brandTeal,
    arc=6pt, boxrule=1.2pt, halign=center, halign title=center,
    coltitle=white, coltext=brandDark, fonttitle=\khatt\bfseries,
    top=1.5mm, bottom=1.5mm, left=2mm, right=2mm, title={#2}, #1}
% Prompt text written to the tool
\newtcolorbox{amr}[1][brandTeal]{enhanced, arc=1mm, boxrule=0pt,
    colback=white, overlay={\fill[#1] (frame.north east) rectangle ([xshift=-2.5pt]frame.south east);},
    left=6pt, right=8pt, top=2pt, bottom=2pt,
    before skip=3pt, after skip=3pt}

%------------------------------------------------
% TikZ (Arabic right-to-left inside nodes)
%------------------------------------------------
\newcommand{\nodecolor}{}
\tikzset{
    every node/.append style={
        execute at begin node={\nodecolor\everypar{\beginR}\beginR},
        execute at end node=\endR},
    lawn/.code={\def\nodecolor{\color{#1}}},
    khatwa/.style={draw=#1, fill=#1!12, very thick, rounded corners=5pt,
        minimum height=0.95cm, text width=1.9cm, align=center,
        font=\khatt\bfseries\small},
    khatwa/.default=brandTeal,
    sahm/.style={-{Stealth[length=2.5mm]}, thick, color=brandMuted},
    rang/.style={draw=#1, fill=#1!12},
    raqm/.style={circle, fill=#1, lawn=white, inner sep=1.5pt,
        font=\small\bfseries},
}

% Horizontal chain of steps, right to left:
% \flow[width]{{a}, {b\\c}, {d}}
\newcommand{\flow}[2][1.9cm]{%
    \begin{tikzpicture}[node distance=5mm]
        \foreach \t [count=\i] in {#2} {
            \ifnum\i=1
                \node[khatwa, text width=#1] (f\i) {\t};
            \else
                \pgfmathtruncatemacro\p{\i-1}
                \node[khatwa, text width=#1, left=of f\p] (f\i) {\t};
                \draw[sahm] (f\p) -- (f\i);
            \fi
        }
    \end{tikzpicture}}

%------------------------------------------------
% Cover (template's 1st slide on the light background)
%------------------------------------------------
\newcommand{\cover}[3]{%
\begin{frame}[plain]
    \begin{tikzpicture}[remember picture, overlay]
        \draw[brandTeal, line width=1.5pt, opacity=0.3]
            (current page.south east) ++(-1,1) circle (4cm);
        \draw[brandTeal, line width=0.8pt, opacity=0.25]
            (current page.north west) ++(2,2) circle (6cm);
        \draw[brandOrange, line width=0.5pt, opacity=0.25]
            (current page.north west) ++(2,2) circle (7cm);
        \node[anchor=north east] at ($(current page.north east)+(-1.0,-0.8)$) {
            \begin{tikzpicture}[scale=0.35]
                \fill[brandTeal] (0,0) -- (0.8,1.2) -- (1.6,0) -- cycle;
                \fill[brandOrange] (0.8,1.2) -- (1.6,2.4) -- (2.4,1.2) -- cycle;
            \end{tikzpicture}
        };
    \end{tikzpicture}

    \vspace{-0.05cm}
    {\hspace*{1.1cm}\khatt\bfseries\color{brandDark}
        المدرسة العليا للأساتذة ببوزريعة ــ ملحقة بومرداس\par}

    \vspace{1.2cm}
    \begin{minipage}{0.85\textwidth}
        {\color{tealDeep}\khatt\bfseries #1} \\[0.3cm]
        {\raggedleft\color{brandDark}\Huge\khatt\bfseries #2\par} \vspace{0.3cm}
        {\color{brandMuted}\large #3}
    \end{minipage}

    \vspace{1.2cm}
    \begin{minipage}{0.85\textwidth}
        \color{brandMuted}\small
        \begin{tabular}{@{}ll}
            \textbf{الأستاذ:} & د. بن عبد الرزاق يوسف \\
            \textbf{الفئة:}   & طلبة قسم اللغة العربية \\
        \end{tabular}
    \end{minipage}
\end{frame}}

\title{التفاعل الفعّال مع الذكاء الاصطناعي}

\begin{document}

\cover{درس الفصل الثالث}
{التفاعل الفعّال مع الذكاء الاصطناعي}
{مبادئ هندسة الأوامر \textenglish{(Prompt Engineering)}}

\begin{frame}{أهداف الدرس}
    \begin{tbox}{\faBullseye\quad أهداف الدرس}
        في نهاية هذا الدرس، تستطيع أن:
        \begin{itemize}
            \item تعرّف الأمر وهندسة الأوامر.
            \item تذكر العناصر الخمسة للأمر الجيد.
            \item تحوّل أمرًا عامًا إلى أمر دقيق خطوة بخطوة.
            \item تستخرج عناصر أي أمر مكتوب.
        \end{itemize}
    \end{tbox}
\end{frame}

\begin{frame}{وضعية للتفكير}
    \small
    كتب طالبان طلبًا في الموضوع نفسه:
    \begin{obox}{\faTimesCircle\quad الطالب الأول}
        أعطني خطة للمراجعة.
    \end{obox}
    \begin{tbox}{\faCheckCircle\quad الطالب الثاني}
        تصرّف كمساعد أكاديمي. أنشئ خطة مراجعة لمادة النحو مدتها 7 أيام،
        بمعدل ساعتين يوميًا. حدّد موضوع كل يوم، وأضف تمرينًا قصيرًا، ونظّم الخطة في جدول.
    \end{tbox}
    الطالب الثاني سيحصل على إجابة أفضل بكثير، لأنه \textbf{أعطى الأداة كل ما تحتاجه} لتفهم طلبه.
\end{frame}

\begin{frame}{ما هندسة الأوامر؟}
    \begin{tbox}{\faLightbulb\quad تعريف}
        \textbf{الأمر} \textenglish{(Prompt)}: النص الذي نكتبه للأداة لنطلب منها شيئًا.\\[3pt]
        \textbf{هندسة الأوامر}: مهارة كتابة أوامر واضحة ومنظمة،
        للحصول على أفضل إجابة ممكنة.
    \end{tbox}
    الأداة لا تعرف ما في ذهنك. هي تعرف فقط ما تكتبه لها،
    لذلك يجب أن يكون الأمر \textbf{واضحًا وكاملًا}.
\end{frame}

\begin{frame}{العناصر الخمسة للأمر الجيد}
    \begin{center}
    \begin{tikzpicture}[
        unsur/.style={draw=brandTeal, fill=brandTeal!12, thick, circle,
            minimum size=1.6cm, align=center, font=\khatt\bfseries}]
        \node[draw=brandOrange, fill=brandOrange!15, very thick, rounded corners=6pt,
            minimum width=2.4cm, minimum height=0.9cm,
            font=\khatt\bfseries\large] (m) {الأمر الجيد};
        \node[unsur] (r) at (90:2.4)   {الدور};
        \node[unsur] (t) at (18:2.7)   {المهمة};
        \node[unsur] (s) at (-54:2.5)  {السياق};
        \node[unsur] (q) at (-126:2.5) {القيود};
        \node[unsur] (f) at (162:2.7)  {التنسيق};
        \foreach \x in {r,t,s,q,f} \draw[thick, brandMuted] (m) -- (\x);
    \end{tikzpicture}
    \end{center}
\end{frame}

\begin{frame}{العناصر الخمسة بالتفصيل}
    \begin{tabularx}{\textwidth}{|>{\centering\arraybackslash}p{2cm}|R|R|}
        \hline
        \ras{العنصر} & \ras{السؤال الذي يجيب عنه} & \ras{مثال} \\ \hline
        \textbf{الدور}    & من تكون الأداة؟ أي خبرة تعتمد؟ & تصرّف كأستاذ لغة عربية. \\ \hline
        \textbf{المهمة}   & ماذا نريد أن تفعل الأداة؟ & اشرح درس الحال. \\ \hline
        \textbf{السياق}   & ما المعلومات التي تساعدها على الفهم؟ & التلاميذ في السنة الثانية متوسط. \\ \hline
        \textbf{القيود}   & ما الشروط التي يجب احترامها؟ & في 150 كلمة، بلغة بسيطة. \\ \hline
        \textbf{التنسيق}  & كيف نريد شكل الإجابة؟ & تعريف، ثم أمثلة، ثم خلاصة. \\ \hline
    \end{tabularx}
\end{frame}

\begin{frame}{بناء الأمر خطوة بخطوة}
    نبدأ بأمر بسيط، ثم نضيف عنصرًا في كل مرة. كل درجة تجعل الأمر أدق:
    \begin{center}
        \flow[2.3cm]{{أمر عام}, {+ السياق}, {+ المهمة\\والقيود}, {+ التنسيق}}
    \end{center}
    \begin{center}
        {\color{tealDeep}\faArrowLeft}\quad \textbf{جودة أفضل للإجابة في كل درجة}
    \end{center}
\end{frame}

\begin{frame}{الدرجات الأربع}
    \small
    \textbf{\color{brandOrange}الدرجة 1: أمر عام}
    \begin{amr}[brandOrange]
        اشرح الحال.
    \end{amr}
    \textbf{\color{tealDeep}الدرجة 2: نضيف السياق}
    \begin{amr}[brandTeal!45]
        أنا أستاذ متربص. تلاميذي في السنة الثانية متوسط. اشرح لي درس الحال.
    \end{amr}
    \textbf{\color{tealDeep}الدرجة 3: نضيف المهمة والقيود}
    \begin{amr}[brandTeal!70]
        \dots اشرح درس الحال في 150 كلمة، بلغة بسيطة، مع مثالين من الحياة اليومية.
    \end{amr}
    \textbf{\color{tealDeep}الدرجة 4: نضيف التنسيق}
    \begin{amr}[brandTeal]
        \dots نظّم الإجابة في ثلاثة أجزاء: تعريف، أمثلة، خلاصة.
    \end{amr}
\end{frame}

\begin{frame}{كيف نقرأ أمرًا ونستخرج عناصره؟}
    \begin{amr}[brandTeal]
        تصرّف كأستاذ جامعي في البلاغة. اشرح مفهوم التشبيه لطالب في السنة الأولى،
        في أقل من 200 كلمة، مع جدول يبيّن أركان التشبيه.
    \end{amr}
    \begin{center}
    \begin{tabularx}{0.75\textwidth}{|R|>{\centering\arraybackslash}p{2.5cm}|}
        \hline
        \ras{الجزء من الأمر} & \ras{العنصر} \\ \hline
        تصرّف كأستاذ جامعي في البلاغة & الدور \\ \hline
        اشرح مفهوم التشبيه            & المهمة \\ \hline
        لطالب في السنة الأولى         & السياق \\ \hline
        في أقل من 200 كلمة            & القيود \\ \hline
        مع جدول يبيّن أركان التشبيه   & التنسيق \\ \hline
    \end{tabularx}
    \end{center}
\end{frame}

\begin{frame}{أخطاء شائعة}
    \begin{obox}{\faExclamationTriangle\quad تجنّب هذه الأخطاء}
        \begin{itemize}
            \item \textbf{أمر قصير جدًا:} «اشرح البلاغة» لا يحدد المستوى ولا الهدف.
            \item \textbf{أمر غامض:} «اكتب شيئًا جميلًا» لا يحدد ما المطلوب.
            \item \textbf{عدة مهام مختلطة:} الأفضل أن نقسم المهام الكبيرة إلى أوامر متتابعة.
            \item \textbf{قبول أول إجابة:} يمكن دائمًا أن نعدّل الأمر ونعيد المحاولة.
        \end{itemize}
    \end{obox}
\end{frame}

\begin{frame}{أوامر مفيدة لأستاذ اللغة العربية}
    \begin{tbox}{\faChalkboardTeacher\quad أمثلة جاهزة}
        \begin{itemize}
            \item اقترح عشر جمل فيها مفعول لأجله، لتلاميذ السنة الثالثة متوسط.
            \item حوّل هذه الأبيات إلى نثر بسيط، دون تغيير المعنى.
            \item أنشئ خمسة أسئلة فهم حول النص الآتي، مع الإجابات.
        \end{itemize}
    \end{tbox}
\end{frame}

\begin{frame}{خلاصة الدرس}
    \begin{tbox}{\faCheckCircle\quad خلاصة الدرس}
        \begin{itemize}
            \item \textbf{الأمر الجيد} يعطي الأداة كل ما تحتاجه لتفهم الطلب.
            \item عناصره الخمسة: \textbf{الدور}، \textbf{المهمة}، \textbf{السياق}،
                  \textbf{القيود}، \textbf{التنسيق}.
            \item نبني الأمر \textbf{تدريجيًا}، ونحسّنه بعد قراءة الإجابة.
        \end{itemize}
    \end{tbox}
\end{frame}

\end{document}
